\documentclass[11pt]{article}

\usepackage{amsmath,amssymb,amsthm,mathtools}

\newtheorem{lemma}{Lemma}

\begin{document}

\begin{lemma}[List-schedule upper bound under pair separation]
Let \(m\ge 2\) be the number of identical parallel machines, and let
\(\alpha>0\). Consider an instance with processing times
\(p_1,\dots,p_n>0\). For a schedule \(S\), define
\[
F_\alpha(S)=\sum_{j=1}^n p_j C_j(\alpha),
\]
where \(C_j(\alpha)\) denotes the load preceding job \(j\) on its machine plus
\(\alpha p_j\).

Let \(W^{LS}_{\alpha,m}\) denote the maximum value of \(F_\alpha(S)\) over all
schedules obtainable by list scheduling over all job orders. Fix a list
schedule attaining \(W^{LS}_{\alpha,m}\). Let \(OPT^R_{\alpha,m}\) denote the
minimum value of \(F_\alpha(S)\) among all feasible schedules satisfying the
following pair-separation restriction: if two jobs are scheduled on the same
machine in the fixed worst list schedule, then they may not be scheduled on
the same machine in the restricted optimum.

Assume that such a restricted optimum exists. Then
\[
\frac{W^{LS}_{\alpha,m}}{OPT^R_{\alpha,m}}
\le
1+\frac{m-1}{2\alpha}.
\]
\end{lemma}

\begin{proof}
For any schedule \(S\), write
\[
S_2=\sum_{j=1}^n p_j^2.
\]
The objective admits the pairwise representation
\[
F_\alpha(S)
=
\alpha S_2
+
\sum_{\{j,k\}: j,k \text{ are on the same machine in } S} p_jp_k.
\]

Indeed, on a machine containing jobs with processing times
\(x_1,\dots,x_q\), the contribution to \(F_\alpha\) is
\[
\sum_{r=1}^q x_r
\left(
\sum_{s<r}x_s+\alpha x_r
\right)
=
\alpha\sum_{r=1}^q x_r^2
+
\sum_{1\le s<r\le q}x_sx_r.
\]
Summing over all machines gives the claimed representation.

Now consider the fixed list schedule attaining \(W^{LS}_{\alpha,m}\). Let its
machines be \(W_1,\dots,W_m\). For each \(i\), let \(q_i=|W_i|\). Since the
restricted optimum must separate every pair of jobs that lies on the same
machine \(W_i\), all jobs in \(W_i\) must be assigned to distinct machines in
the restricted optimum. Since there are only \(m\) machines, feasibility
implies
\[
q_i\le m
\qquad
\text{for every } i.
\]

For a fixed list-schedule machine \(W_i\), let
\[
P_i=\sum_{j\in W_i}p_j,
\qquad
S_i=\sum_{j\in W_i}p_j^2.
\]
The pairwise contribution of \(W_i\) is
\[
\sum_{\{j,k\}\subseteq W_i}p_jp_k
=
\frac12\left(P_i^2-S_i\right).
\]
By Cauchy--Schwarz,
\[
P_i^2
=
\left(\sum_{j\in W_i}p_j\right)^2
\le
q_i \sum_{j\in W_i}p_j^2
=
q_i S_i.
\]
Therefore,
\[
\sum_{\{j,k\}\subseteq W_i}p_jp_k
\le
\frac12(q_iS_i-S_i)
=
\frac{q_i-1}{2}S_i.
\]
Since \(q_i\le m\), this gives
\[
\sum_{\{j,k\}\subseteq W_i}p_jp_k
\le
\frac{m-1}{2}S_i.
\]

Summing over all machines of the fixed worst list schedule yields
\[
\sum_{\{j,k\}: j,k \text{ are on the same machine in the list schedule}}
p_jp_k
\le
\frac{m-1}{2}\sum_{i=1}^m S_i
=
\frac{m-1}{2}S_2.
\]
Hence
\[
W^{LS}_{\alpha,m}
\le
\alpha S_2+\frac{m-1}{2}S_2.
\]

On the other hand, for every feasible restricted schedule \(S\),
\[
F_\alpha(S)
=
\alpha S_2
+
\sum_{\{j,k\}: j,k \text{ are on the same machine in }S}p_jp_k
\ge
\alpha S_2.
\]
Therefore,
\[
OPT^R_{\alpha,m}\ge \alpha S_2.
\]

Combining the two inequalities gives
\[
\frac{W^{LS}_{\alpha,m}}{OPT^R_{\alpha,m}}
\le
\frac{\alpha S_2+\frac{m-1}{2}S_2}{\alpha S_2}
=
1+\frac{m-1}{2\alpha}.
\]
This proves the claim.
\end{proof}

\end{document}